% Generated draft: requires mathematical and source review.
\documentclass[11pt]{article}
\usepackage[margin=25mm]{geometry}
\usepackage{fontspec}
\setmainfont{Latin Modern Roman}
\usepackage{amsmath,amssymb,mathtools,unicode-math}
\setmathfont{Latin Modern Math}
\usepackage{enumitem,needspace,xurl,hyperref}
\setlength{\parindent}{0pt}
\setlength{\parskip}{5pt}
\setlength{\emergencystretch}{3em}
\newcommand{\N}{\mathbb{N}}
\newcommand{\Z}{\mathbb{Z}}
\newcommand{\Q}{\mathbb{Q}}
\newcommand{\R}{\mathbb{R}}
\newcommand{\C}{\mathbb{C}}
\newcommand{\F}{\mathbb{F}}
\newcommand{\A}{\mathbb{A}}
\newcommand{\PP}{\mathbb{P}}
\newcommand{\E}{\mathbb{E}}
\newcommand{\eps}{\varepsilon}
\newcommand{\dd}{\,\mathrm{d}}
\begin{document}
\section*{UnsolvedMath ID: 30006228}
\textbf{Source problem code:} catalogue entry
\subsection{Definitions and mathematical statement}
Let \(N\geq 1\) be an integer and let \(|x|\) denote the Euclidean norm on \(\R^N\). For an essentially bounded potential \(V\), write
\[
H=-\Delta+V,\qquad \Delta=\sum_{j=1}^{N}\frac{\partial^2}{\partial x_j^2},\qquad
\|V\|_{L^\infty(\R^N)}=\operatorname*{ess\,sup}_{x\in\R^N}|V(x)|.
\]
The reconstructed source question asks whether, without a positivity assumption on \(H\), every solution of
\[
Hu=0\quad\text{in }\R^N,\qquad \|V\|_{L^\infty(\R^N)}\leq 1,
\]
that decays at least as fast as \(e^{-k|x|}\) for some \(k>1\) must satisfy \(u\equiv0\).

Under the usual uniform big-\(O\) interpretation of the decay phrase, its hypothesis is
\[
\exists k>1\;\exists C>0\;\exists R>0\;\forall x\in\R^N:
\quad |x|\geq R\ \Longrightarrow\ |u(x)|\leq C e^{-k|x|}.
\]
This interpretation is explicit here but is not specified by the supplied wording. For a solution defined only almost everywhere, the corresponding interpretation uses an almost-everywhere inequality outside the ball. The domain is all of \(\R^N\); no finite boundary or boundary data are supplied.

\textbf{Unresolved formulation qualifications.} The reconstructed statement does not specify whether \(V\) and \(u\) are real-valued or complex-valued, or the regularity and solution concept for \(u\). The supplied literature assessment refers to the real-valued Landis conjecture, suggesting a real-valued intended setting, but this does not unambiguously settle both scalar domains. Thus the displayed implication is a source-faithful question requiring these choices before it becomes a fully specified conjecture; it is not an assertion for every possible scalar domain or solution class.

\textbf{Qualified background from the supplied excerpt.} The excerpt separately describes a partial affirmative result under \(H\geq0\) in \(\R^N\), where this notation is defined there by the existence of a positive supersolution of \(H\phi=0\). Its listed hypotheses are
\[
N\geq1,\qquad V\leq1,\qquad V\in L^q_{\mathrm{loc}}(\R^N),\qquad q>N/2,
\]
with a solution apparently belonging to \(W^{1,2}_{\mathrm{loc}}(\R^N)\). Here \(L^q_{\mathrm{loc}}\) and \(W^{1,2}_{\mathrm{loc}}\) denote local Lebesgue and local first-order Sobolev spaces, respectively. The excerpt also appears to list the growth bounds
\[
|u(x)|=
\begin{cases}
O(1),&N=1,\\
O\!\left(|x|^{(2-N)/2}\right),&N\geq2,
\end{cases}
\qquad |x|\to\infty.
\]
These lines are damaged and incomplete in the supplied extraction. They do not provide the complete sharp decay criterion and must not be presented as sufficient hypotheses for a vanishing theorem. In particular, the background condition \(V\leq1\) is a one-sided bound and is distinct from the reconstructed question's condition \(\|V\|_\infty\leq1\). The source describes its positivity-dependent criterion as weaker than Landis's proposed decay condition; no missing criterion is reconstructed here.
\subsection{Short English statement}
The supplied Landis question asks whether a solution on all of \(\R^N\), with \(\|V\|_\infty\leq1\), must vanish if it decays uniformly like \(O(e^{-k|x|})\) for some \(k>1\), without assuming positivity of \(-\Delta+V\). The scalar domains, solution class, and exact decay convention require clarification. The supplied positivity-dependent partial result has an incomplete extracted statement.
\subsection{Sources}
\begin{itemize}
\item Supplied local UnsolvedMath registry record, permanent ID \(30006228\), title ``Landis Unique-Continuation Conjecture,'' exact source problem code \texttt{catalogue entry}, in the Oberwolfach Reports Open Problems collection. The record attributes the conjectural claim to Landis; no individual proposer is identified.
\item Supplied report citation: \emph{Mini-Workshop: Hardy Inequalities in Discrete and Continuum Settings} (2025), \emph{Oberwolfach Reports}, pp.\ 579--614, DOI 10.4171/owr/2025/14. The extracted passage includes the printed page number 587. \url{https://doi.org/10.4171/owr/2025/14}.
\item Supplied literature reference: Das--Pinchover, ``Landis via criticality theory,'' described by the record as 2024 work. \url{https://arxiv.org/abs/2405.11695}.
\item Supplied literature reference: Rossi, ``Landis with sharp decay.'' \url{https://doi.org/10.1512/iumj.2021.70.8218}.
\item Supplied literature reference: ``Oberwolfach Report 14/2025.'' \url{https://doi.org/10.4171/owr/2025/14}.
\item Literature assessment date: \texttt{2026-08-22}. This is supplied dated curation, not an independently verified assessment. It labels the problem partially solved and states that Das--Pinchover prove a sharp decay criterion for nonnegative Schr\"odinger operators under positivity and related hypotheses, while the general real-valued Landis conjecture remains open outside such regimes. The supplied background explicitly requests independent verification of time-sensitive status and cited claims.
\item The registry labels the statement \texttt{corrected\_verified} and reports that it was reconstructed and checked against the official Oberwolfach report. These are supplied provenance claims; no independent source consultation or mathematical verification is asserted here.
\end{itemize}
\end{document}
